% =====================================================================
% Ensaio 1 — documento de qualificação
%
% ⚠️ Alvo declarado no roteiro (docs/ars/04-roteiro-qualificacao.md):
% "documento de projeto defensável. Identificação fechada, dose verificada,
% descritivas preliminares, ameaças mapeadas. NÃO resultado estimado."
%
% Portanto a seção de identificação está COMPLETA sem dado, e as seções de
% resultado estão ausentes de propósito — não por atraso.
%
% Compilar:  latexmk -pdf -interaction=nonstopmode paper/main.tex
% =====================================================================
\documentclass[12pt,a4paper]{article}

\usepackage[T1]{fontenc}
\usepackage[utf8]{inputenc}
\usepackage[brazil]{babel}
\usepackage{amsmath,amssymb}
\usepackage{booktabs}
\usepackage{graphicx}
\usepackage[margin=3cm]{geometry}
\usepackage[colorlinks=true,linkcolor=black,citecolor=black]{hyperref}
\usepackage{natbib}
\usepackage{setspace}

\onehalfspacing
\graphicspath{{figures/}}

\title{Proibição da pulverização aérea de agrotóxicos no Ceará:\\
       avaliação de impacto sobre desfechos perinatais}
\author{Felipe Carvalho Souza Santos\\
        \small PPGEco/UFES}
\date{\today}

\begin{document}
\maketitle

% 1. Introdução                        <- a escrever por último (fórmula de Head)
% =====================================================================
% Ensaio 1 — Seção 2: Background
%
% Bellemare dispensa a seção de Background em artigo de journal, mas a considera
% útil "quando os detalhes de uma legislação precisam ser tidos em mente para
% avaliar os efeitos dessa legislação sobre algum desfecho". É exatamente o caso.
%
% ⚠️ CORREÇÃO SUBSTANTIVA em relação ao esboço do projeto: o framework é centrado
% em GLIFOSATO. A química documentada na Chapada do Apodi NÃO é essa. Ver §2.3.
% =====================================================================

\section{Background}
\label{sec:background}

\subsection{A Lei Zé Maria do Tomé e sua cronologia}
\label{sub:lei}

Em 08 de janeiro de 2019 o governador do Ceará sancionou a Lei Estadual nº
16.820, publicada no Diário Oficial e vigente a partir de 09 de janeiro do mesmo
ano.\footnote{A distinção entre sanção e vigência é irrelevante em painel mensal
--- ambas caem em janeiro de 2019 ---, mas é registrada aqui porque o texto
deve dizer 09/01/2019 ao falar de vigência e 08/01/2019 ao falar de sanção.} A lei
acrescenta o art. 28-B à Lei Estadual nº 12.228/1993, norma-mãe sobre agrotóxicos
no estado, com a seguinte redação:

\begin{quote}
\textbf{Art. 28-B.} É vedada a pulverização aérea de agrotóxicos na agricultura
no Estado do Ceará.

\textbf{§ 1º} A infração ao art. 1º sujeita o infrator ao pagamento de multa de
15 mil (quinze mil) UFIRCEs.

\textbf{§ 2º} Fica proibida a incorporação de mecanismos de controle vetorial por
meio de dispersão por aeronave em todo o Estado do Ceará, inclusive para os casos
de doenças causadas por vírus.
\end{quote}

Dois traços da redação organizam todo o desenho empírico. Primeiro, a proibição é
de \textbf{método}: nenhum princípio ativo é nomeado. Segundo, ela não se
restringe à agricultura --- o §2º alcança a dispersão aérea sanitária, o que
significa que municípios sem agricultura pulverizada, mas com controle aéreo de
endemias, também são alcançados pela norma.

O Ceará foi o primeiro estado brasileiro a adotar a medida. A lei leva o nome de
José Maria Filho, conhecido como Zé Maria do Tomé, líder comunitário da Chapada
do Apodi assassinado em 2010 após denunciar publicamente a pulverização aérea na
região. O dado não é ornamental: ele situa a proibição como resposta a um
conflito distributivo, e não como ajuste técnico --- ponto ao qual a
§\ref{sec:teoria} retorna.

\begin{table}[htbp]
\centering
\caption{Cronologia legal do tratamento}
\label{tab:cronologia}
\small
\begin{tabular}{lp{10.5cm}}
\hline
Data & Evento \\
\hline
09/12/1993 & Lei Estadual nº 12.228/1993 --- norma-mãe sobre agrotóxicos no Ceará. \\
2009 & Limoeiro do Norte proíbe a pulverização aérea por lei municipal, ao amparo do art. 29 da lei estadual, que autoriza legislação supletiva. \\
2010 & Assassinato de Zé Maria do Tomé. \\
\textbf{24/02/2015} & \textbf{Apresentação do Projeto de Lei nº 18/2015.} Primeiro sinal público --- \emph{marco de notícia}. \\
\textbf{18/12/2018} & \textbf{Aprovação por unanimidade na Assembleia Legislativa}, após quatro anos de tramitação --- \emph{marco de certeza}. \\
\textbf{09/01/2019} & \textbf{Vigência da Lei nº 16.820/2019} --- \emph{marco de obrigação}. \\
mai/2023 & O Supremo Tribunal Federal julga a lei constitucional (ADI 6137, unânime). \\
19/12/2024 & Lei Estadual nº 19.135/2024 dá nova redação ao art. 28-B, abrindo exceção para aeronaves remotamente pilotadas. \\
2025 & Nova contestação no STF (ADI 7794), pendente. \\
\hline
\end{tabular}
\end{table}

\paragraph{Três marcos, não um.} A Tabela~\ref{tab:cronologia} distingue
deliberadamente notícia, certeza e obrigação. A distinção tem consequência
empírica direta: a janela de medida da dose pré-ban começa \emph{depois} do
primeiro marco, de modo que a variável de tratamento pode já incorporar resposta
antecipatória. A §\ref{sub:dose} desenvolve o ponto.

\paragraph{A reversão de 2024, e por que ela fecha a janela.} A Lei nº
19.135/2024 mantém a vedação, \emph{"salvo se realizada por meio de Aeronaves
Remotamente Pilotadas --- ARPs, Veículo Aéreo Não Tripulado --- VANT ou Drones"},
sob condições de altura, velocidade de vento e distância de escolas, hospitais e
áreas protegidas. A nova redação \textbf{não reproduz} o §2º de 2019: a proibição
de dispersão aérea para controle vetorial cai. Do ponto de vista de instrumento
de política, o Ceará não revogou a proibição --- \textbf{substituiu o
instrumento}, trocando cota zero por regulação de desempenho com zona-tampão.
Para o Ensaio 1, o efeito é fixar o fim da janela de análise em 19/12/2024, de
modo a preservar a cota zero como tratamento.

\subsection{A Chapada do Apodi e a fruticultura irrigada}
\label{sub:chapada}

A Chapada do Apodi, no baixo Jaguaribe, concentra o agronegócio de fruticultura
irrigada do Ceará --- notadamente banana e melão, em perímetros de irrigação
sobre solo cárstico. É a região à qual a lei responde, e é onde a exposição
populacional à deriva foi documentada antes da proibição.

Duas características locais importam ao desenho. A primeira é que a cultura
dominante é \textbf{fruta}, não grão: a intensidade de pulverização está ligada a
controle fitossanitário de ciclo curto, com aplicações repetidas ao longo da
safra. A segunda é hidrogeológica: o Aquífero Jandaíra é cárstico, e a
vulnerabilidade à percolação é heterogênea no território --- o que abre variação
explorável no canal-água, e simultaneamente adverte contra tratar todos os
municípios como igualmente permeáveis.

\subsection{O químico, e uma correção que o desenho exige}
\label{sub:quimico}

\begin{quote}
\textbf{Advertência ao leitor familiarizado com a literatura de agrotóxicos e
saúde.} Boa parte dessa literatura --- e, no Brasil, boa parte dos desenhos de
avaliação --- é construída sobre \textbf{glifosato} aplicado a culturas
transgênicas tolerantes a herbicida. \textbf{Não é o caso aqui}, e supor que seja
transferiria mecanismo errado.
\end{quote}

A justificativa do Projeto de Lei nº 18/2015 reproduz tabela de análises de água
oriunda do Dossiê ABRASCO, referente a 23 pontos de coleta na Chapada do Apodi.
A composição encontrada é a seguinte:

\begin{table}[htbp]
\centering
\caption{Princípios ativos detectados em água, Chapada do Apodi}
\label{tab:abrasco}
\begin{tabular}{llc}
\hline
Princípio ativo & Classe & Amostras positivas \\
\hline
Procimidona & Fungicida & 23 de 23 \\
Carbaril & Inseticida (carbamato) & 23 de 23 \\
Carbofurano & Inseticida (carbamato) & 18 de 23 \\
Fenitrotiona & Inseticida (organofosforado) & 16 de 23 \\
Glifosato & Herbicida & 4 de 23 \\
\hline
\end{tabular}

\vspace{0.5em}
\begin{minipage}{0.85\textwidth}
\footnotesize \emph{Fonte:} tabela reproduzida na justificativa do PL nº
18/2015, atribuída ao Dossiê ABRASCO; coleta referente a 2009. \emph{Nota:} a
verificação da fonte primária permanece pendente.
\end{minipage}
\end{table}

O perfil é de \textbf{fungicida e inseticida sobre fruticultura irrigada}, não de
herbicida sobre grão transgênico. Três consequências metodológicas decorrem
disso, e todas são levadas adiante nas seções seguintes:

\begin{enumerate}
  \item \textbf{Os desenhos de referência não transferem quimicamente.} Estudos
    construídos sobre a difusão de soja tolerante a herbicida transferem na
    \emph{forma} --- a lógica de instrumentar exposição por aptidão agroclimática
    ---, não no \emph{mecanismo}. A molécula, a via de exposição e a cultura são
    outras.

  \item \textbf{A receita de instrumentação precisa ser adaptada.} A construção
    do instrumento de aptidão agroclimática, na literatura de referência, toma o
    máximo entre culturas \textbf{transgênicas}. Banana não é transgênica. O
    método transfere; a lista de culturas, não. Verificar se o catálogo
    agroclimático cobre banana e melão é, por isso, pré-requisito e não detalhe.

  \item \textbf{Carbamatos e organofosforados têm assinatura toxicológica
    distinta.} São inibidores de colinesterase, com quadro agudo característico
    --- o que torna a série de internação por intoxicação aguda um canal
    informativo, e não apenas um desfecho secundário.
\end{enumerate}

\subsection{O mecanismo: deriva}
\label{sub:deriva}

O que distingue a aplicação aérea da terrestre é a \textbf{deriva}: parte da
calda não atinge o alvo e é carregada pelo vento, dispersando-se sobre
residências, escolas, corpos d'água e comunidades a favor do vento. A exposição
resultante é \textbf{difusa e populacional}, não ocupacional e pontual --- e é
essa característica que faz da proibição um instrumento de saúde pública, e não
apenas de segurança do trabalho.

Dois canais operam em paralelo, e distingui-los é contribuição de mecanismo:

\begin{description}
  \item[Canal-ar.] Deriva direta sobre populações situadas a favor do vento em
    relação à área de aplicação. A direção do vento é, aqui, fonte potencial de
    identificação: populações a favor e contra o vento diferem em exposição sem
    diferir sistematicamente em outros atributos.

  \item[Canal-água.] Escoamento e percolação para bacias hidrográficas, com
    chegada à água de consumo. A estrutura montante/jusante dentro de bacia
    oferece comparação interna, e a heterogeneidade cárstica do Jandaíra modula a
    intensidade da percolação.
\end{description}

\paragraph{Um terceiro canal, de sinal oposto.} Proibido o avião, a aplicação não
desaparece: migra para trator e pulverizador costal. Isso \emph{afasta} o produto
da população --- menos deriva --- e simultaneamente o \emph{aproxima} do
aplicador. A proibição pode, portanto, reduzir exposição populacional e
\textbf{aumentar} intoxicação aguda ocupacional. O canal de substituição é
medido separadamente, e sua direção esperada é oposta à do desfecho principal.

\subsection{Fiscalização}
\label{sub:enforcement}

A fiscalização do art. 28-B cabe, pelos arts. 15 e 30 da Lei nº 12.228/1993, à
Superintendência Estadual do Meio Ambiente, à Secretaria da Agricultura e à
Secretaria da Saúde. O art. 8º da mesma lei determina o registro, junto ao órgão
ambiental estadual, das empresas prestadoras de serviço que utilizam agrotóxicos
--- registro que, se público e com série histórica municipal, constitui a única
fonte estadual capaz de identificar quem operava aplicação aérea antes da
proibição.

A extensão efetiva da fiscalização não é conhecida. Trata-se de limitação
declarada, com consequência precisa sobre o estimando: se a proibição não foi
fiscalizada, a dose pré-ban mede \emph{intenção} de aplicação, não exposição
efetivamente removida, e o parâmetro estimado passa a ser um efeito de intenção
de tratar. A §\ref{sub:ameacas} trata do ponto, e a §\ref{sub:deriva} indica a
via independente pela qual ele pode ser parcialmente contornado.

% =====================================================================
% Ensaio 1 — Seção 3: Referencial teórico
%
% Vai do geral ao específico e TERMINA numa predição testável — que é o que
% amarra esta seção à §4 (identificação) e planta o Ensaio 2.
%
% ⚠️ REFERÊNCIAS NÃO VERIFICADAS. As citações abaixo vêm de
% docs/framework-dissertacao.md, documento do pesquisador. Nenhuma foi conferida
% contra DOI/Crossref — as bases estão inacessíveis (ver docs/lacunas-de-dados.md
% §1, classe E). Conferir antes da versão final; o CLAUDE.md é explícito.
% =====================================================================

\section{Referencial teórico}
\label{sec:teoria}

\subsection{A deriva como externalidade negativa}
\label{sub:externalidade}

O agrotóxico que deriva sobre o vizinho é o caso de manual de externalidade
negativa: um custo imposto a terceiros que não transita pelo sistema de preços. O
produtor que contrata a aplicação aérea internaliza o custo do serviço e do
produto, não o custo em saúde suportado pela população a favor do vento. O
resultado privado é aplicação acima do nível socialmente eficiente.

As duas respostas canônicas divergem no instrumento e, mais fundo, no que
consideram o objeto do problema. A resposta pigouviana internaliza o custo por
tributo igual ao dano marginal, fazendo o mercado precificar o que antes não
precificava. A resposta coaseana observa que, com direitos de propriedade bem
definidos e custos de transação baixos, as partes negociam a solução eficiente
por conta própria --- e desloca a pergunta para \emph{quem detém o direito}: o de
pulverizar, ou o de não ser pulverizado.

\paragraph{Por que a barganha coaseana não se aplica aqui.} A deriva atinge
população rural difusa, numerosa e não organizada, sobre território extenso, com
dano probabilístico, latente e de difícil atribuição individual. Os custos de
transação são, por construção, proibitivos. Isso não é detalhe de aplicação: é o
que \emph{justifica} a intervenção regulatória em vez da negociação privada, e é
também o que desloca a discussão para a escolha de instrumento --- objeto do
Ensaio 2.

\subsection{Preço ou quantidade}
\label{sub:weitzman}

O núcleo teórico do Ensaio 2 é a comparação de instrumentos sob incerteza. Sob
incerteza a respeito dos custos de controle, um instrumento de preço --- taxa
sobre o insumo --- e um de quantidade --- cota, ou proibição --- \textbf{não são
equivalentes em bem-estar esperado}. Qual deles domina depende das
\textbf{inclinações relativas} das curvas de benefício e custo marginais.

A regra de bolso: se o benefício marginal do controle é mais inclinado que o
custo marginal --- danos que disparam rapidamente ao ultrapassar um limiar,
plausível em desfechos de saúde ---, o instrumento de \textbf{quantidade} domina.
Se o custo marginal é mais inclinado, o instrumento de \textbf{preço} domina. Uma
proibição é o caso-limite de cota fixada em zero.

\begin{quote}
\textbf{A ligação que amarra os dois ensaios.} Weitzman decide pela
\textbf{inclinação}, não pelo nível. É por isso que o Ensaio 1 tem como
parâmetro-alvo a \emph{curva} dose-resposta e não o efeito num ponto: o Ensaio 2
consome a derivada. Um Ensaio 1 que entregasse apenas o efeito médio no nível
tratado deixaria o Ensaio 2 sem âncora empírica --- e a escolha entre proibir e
taxar voltaria a ser argumentada por plausibilidade, que é exatamente o que a
dissertação se propõe a substituir.
\end{quote}

\paragraph{A camada de realismo: poluição difusa.} Instrumentos baseados em
emissão observada pressupõem que se sabe quem emitiu quanto. A deriva não tem
emissor identificável --- é poluição de fonte difusa. Sob heterogeneidade entre
produtores, a literatura registra que a taxa tende a ser mais eficiente enquanto
a regulação por quantidade gera maior lucro ao setor regulado --- resultado com
consequências distributivas, e não apenas de eficiência. A observação importa
porque a proibição cearense foi contestada judicialmente pelo lobby do setor, e
não por consumidores.

\subsection{A regulação de agrotóxicos em específico}
\label{sub:regulacao}

A literatura específica sobre regulação de pesticidas acrescenta duas coisas que
o instrumental geral não dá. A primeira é a análise marginal dos custos de
bem-estar \emph{e dos efeitos distributivos} da regulação --- quem paga a conta
da restrição, e em que proporção. A segunda é a comparação empírica de
instrumentos, em que há registro de cenários nos quais cotas superam taxas ---
isto é, o resultado de Weitzman aplicado, e não apenas invocado.

\paragraph{Uma advertência que a §\ref{sub:quimico} torna necessária.} A
literatura econômica sobre regulação de agrotóxicos é, majoritariamente,
construída sobre herbicidas em culturas de grão, frequentemente transgênicas.
O caso cearense é de fungicidas e inseticidas sobre fruticultura irrigada. A
estrutura de custos de substituição, a elasticidade de resposta do produtor e o
perfil toxicológico são outros. As predições qualitativas do referencial se
mantêm; as magnitudes emprestadas, não.

\subsection{O que o instrumento pressupõe, conta e exclui}
\label{sub:normativo}

Escolher entre proibir, taxar e regular desempenho é problema de desenho de
instrumento. Tratá-lo apenas como exercício de maximização de bem-estar esperado,
porém, deixa fora do quadro uma camada que este caso torna difícil de ignorar.

Precificar uma externalidade de saúde --- a via da taxa --- e proibi-la --- a via
da cota zero --- não são escolhas técnicas equivalentes com sinais trocados.
Elas diferem no que pressupõem: a primeira supõe que dano à saúde é comensurável
com dinheiro e que a compensação é a forma adequada de resposta; a segunda supõe
que há um direito a não ser exposto, anterior ao cálculo. Diferem também no que
contam: um cálculo welfarista puro registra o custo agregado da restrição e o
benefício agregado em saúde, e não registra quem, especificamente, absorve o
resíduo quando o instrumento falha.

A lei em análise leva o nome de um líder rural assassinado por se opor à
pulverização. Uma avaliação que se limite a estimar o efeito médio sobre peso ao
nascer não está errada --- é, aliás, o que esta dissertação faz, e faz com o rigor
que o método exige. Mas ela é incompleta se não registrar que a escolha de
instrumento é também uma posição sobre de quem é o direito. Esta seção declara a
camada; a §\ref{sec:identificacao} não a usa como argumento empírico, e a
conclusão retorna a ela na discussão de política.

\subsection{A predição testável}
\label{sub:predicao}

O referencial acima produz predições, e é por elas que a §\ref{sec:identificacao}
é avaliada.

\begin{description}
  \item[Predição principal.] Se a pulverização aérea gera externalidade de saúde
    por deriva, sua proibição reduz a exposição populacional, e essa redução é
    \textbf{crescente na intensidade pré-ban} de aplicação. Municípios que mais
    aplicavam devem apresentar a maior melhora em desfechos perinatais. É a
    predição que a curva dose-resposta testa.

  \item[Predição de canal --- água.] Se parte da exposição transita pelo canal
    hídrico, a melhora deve ser maior a jusante do que a montante dentro da mesma
    bacia, para dose pré-ban comparável.

  \item[Predição de canal --- ar.] Se parte transita pelo canal atmosférico, a
    melhora deve ser maior a favor do vento dominante do que contra.

  \item[Predição de substituição, de sinal oposto.] Se a aplicação migra do
    avião para o solo, a intoxicação aguda \textbf{ocupacional} deve
    \textbf{aumentar} com a dose pré-ban, ainda que a exposição populacional
    caia.

  \item[Predição nula, que funciona como falsificação.] Como a proibição alcança
    o \emph{método} e não a \emph{molécula}, a disponibilidade do produto não se
    altera. Logo, a intoxicação \textbf{autoprovocada} --- que responde à
    disponibilidade, não ao modo de aplicação --- \textbf{não} deve se mover com
    a dose. As duas séries partilham população, sistema de registro e municípios;
    movimento conjunto não sustenta a hipótese de substituição, e sim algum fator
    que desloca internação em geral.

  \item[Predição de custo.] A proibição impõe custo ao produtor --- substituição
    de tecnologia de aplicação, possível perda de produtividade. É o lado do
    trade-off que o Ensaio 2 precifica, e cuja mensuração fica na agenda.
\end{description}

\paragraph{O que torna o conjunto informativo.} As predições não apontam todas na
mesma direção. Duas delas --- substituição e falsificação --- têm sinais
esperados que \emph{contradizem} a leitura ingênua de que "a proibição melhora
tudo". Um desenho no qual todo resultado possível confirmaria a hipótese não
testaria hipótese alguma; a §\ref{sub:poder} formaliza esse critério ao declarar,
antes da estimação, que resultado constituiria evidência contrária.

% =====================================================================
% Ensaio 1 — Seção 4: Estratégia empírica e identificação
%
% É a seção que a qualificação defende. O alvo declarado no roteiro é
% "identificação fechada, ameaças mapeadas, NÃO resultado estimado" — então esta
% seção está completa mesmo sem dado, e é assim que ela deve ser lida.
%
% Escrita em lógica de forward engineering, como o CLAUDE.md pede:
% parâmetro-alvo -> hipóteses de identificação -> estimação. Não o contrário.
%
% ⚠️ Números de resultado NÃO entram aqui. Onde há número, ele é de DESENHO
% (poder, efeito mínimo detectável, aritmética de exposição) e está marcado.
% =====================================================================

\section{Estratégia empírica e identificação}
\label{sec:identificacao}

\subsection{O parâmetro-alvo, e por que ele custa caro}
\label{sub:estimando}

A Lei Estadual nº 16.820/2019 entrou em vigor em 09 de janeiro de 2019 e vale
para o Ceará inteiro de uma só vez. Não há, portanto, variação de \emph{timing}
de adoção dentro do estado: a máquina de diferenças-em-diferenças escalonado não
se aplica a este evento. A identificação vem de outra fonte de variação --- a
\textbf{intensidade pré-ban de exposição à pulverização aérea} entre os 184
municípios cearenses ---, e o tratamento é, por construção, \textbf{contínuo}.

O estimador é o de Callaway, Goodman-Bacon e Sant'Anna para tratamento contínuo.
Ele separa dois parâmetros que a literatura frequentemente confunde, e a
distinção organiza toda esta seção:

\begin{itemize}
  \item $ATT(d\,|\,d)$ --- o efeito médio, entre os municípios que receberam
    dose $d$, de ter recebido $d$ em vez de zero. É um efeito \textbf{no nível}
    da dose.
  \item $ATE(d)$ e sua derivada $ACR(d)$ --- a \textbf{curva} dose-resposta
    causal: o efeito que $d$ teria sobre um município tomado ao acaso, e a taxa
    a que esse efeito varia com a dose.
\end{itemize}

Os dois exigem hipóteses diferentes. $ATT(d\,|\,d)$ sai sob paralelismo de
tendências no resultado \emph{não tratado} --- a hipótese usual. A curva exige
\textbf{tendências paralelas fortes} (\emph{strong parallel trends}), que
restringe trajetórias sob doses que o município \textbf{não recebeu}. Essa
hipótese é, por natureza, \textbf{não testável}: nenhum período pré-tratamento
informa o que teria acontecido a um município sob uma dose alternativa.

\begin{quote}
\textbf{Registro de honestidade.} O teste de tendências pré-tratamento --- o
\emph{event study} de \emph{leads} --- \textbf{não} testa a hipótese forte. Ele
fala da hipótese usual. Apresentá-lo como validação da curva seria erro de
leitura, e é um erro comum o bastante para merecer estar escrito.
\end{quote}

Há ainda uma consequência econômica. Sob as duas hipóteses conjuntamente,
$ATT(d\,|\,d) = ATE(d)$, o que equivale a excluir \emph{selection-on-gains}:
municípios não podem ter escolhido sua intensidade agrícola antecipando ganhos
diferenciados. Isso é uma afirmação forte sobre decisão de plantio, não um
tecnicismo. Ela é declarada, não assumida em silêncio.

\paragraph{A escolha, e o que a motiva.} O resultado principal é a
\textbf{curva}, acompanhada dos limites de identificação parcial sob hipótese de
direção do viés. A razão não é econométrica: o Ensaio 2 compara instrumentos de
política sob a lógica de Weitzman, que decide pela \textbf{inclinação relativa}
do dano marginal. Ou seja, o segundo ensaio precisa de $ACR(d)$ --- a derivada.
Recuar para $ATT(d\,|\,d)$ tornaria o Ensaio 1 mais barato de defender e deixaria
o Ensaio 2 sem âncora quantitativa. É decisão de arquitetura da dissertação, e
está registrada como tal.

\subsection{A variável de tratamento, e o que ela não mede}
\label{sub:dose}

A dose é a intensidade agrícola municipal pré-ban da cultura-âncora, medida pela
área plantada média de 2015 a 2018 na Produção Agrícola Municipal do IBGE.
Utiliza-se área \textbf{plantada}, e não colhida: o período 2012--2017 foi de
seca severa no semiárido cearense, e a área colhida responde à quebra de safra
--- isto é, é endógena ao clima, que é justamente o canal que o instrumento de
aptidão agroclimática busca isolar. Área plantada mede intenção de plantio.

Três limitações de medida acompanham essa escolha, e nenhuma delas se resolve
com mais dado:

\begin{enumerate}
  \item \textbf{A proibição é de método, não de molécula.} O art. 28-B veda a
    pulverização aérea; não nomeia princípio ativo algum. Área plantada é proxy
    de intensidade agrícola, não de aplicação aérea efetiva, e culturas diferem
    em quanto de fato recebem avião. O viés de medida daí resultante é, em geral,
    atenuante.

  \item \textbf{A química local não é a da literatura de referência.} A tabela
    de análises de água reproduzida na justificativa do Projeto de Lei nº
    18/2015, oriunda do Dossiê ABRASCO, registra em 23 pontos de coleta na
    Chapada do Apodi: procimidona e carbaril em 23 de 23 amostras, carbofurano em
    18, fenitrotiona em 16 --- e glifosato em 4. O perfil é de fungicida e
    inseticida sobre fruticultura irrigada, não de herbicida sobre grão
    transgênico. A consequência é metodológica: os desenhos construídos sobre
    soja tolerante a herbicida e glifosato não transferem quimicamente, ainda que
    transfiram na forma.

  \item \textbf{A janela de medida da dose começa depois da notícia do ban.} O
    Projeto de Lei foi apresentado em 24 de fevereiro de 2015. A janela
    2015--2018, portanto, é integralmente posterior ao primeiro sinal público de
    que a proibição viria. Se produtores responderam à expectativa, a dose medida
    já incorpora parte do ajuste --- e isso morde a \emph{variável de tratamento},
    não apenas o desfecho. Esta é a limitação mais séria das três.
\end{enumerate}

\paragraph{Três marcos, não um.} A cronologia legal distingue notícia
(24/02/2015, apresentação do projeto), certeza (18/12/2018, aprovação unânime na
Assembleia) e obrigação (09/01/2019, vigência). O desenho os trata separadamente:
os \emph{leads} do \emph{event study} alcançam 2018 precisamente para que o
segundo marco seja visível.

\subsection{Quando o tratamento alcança cada coorte}
\label{sub:exposicao}

Um desfecho perinatal não é contemporâneo à exposição. Um nascimento ocorrido em
março de 2019 foi gestado entre junho de 2018 e março de 2019 --- dois terços da
gestação \textbf{antes} da vigência da lei. Marcar como tratada toda coorte
nascida a partir de janeiro de 2019 dilui o efeito com gestação não exposta, e a
diluição tem direção conhecida.

O painel constrói, para cada célula município--mês, a fração da janela
gestacional posterior ao ban. O tratamento não liga de uma vez: sobe ao longo de
aproximadamente nove meses, e apenas as coortes nascidas a partir de outubro de
2019 têm gestação integralmente pós-ban.

\begin{table}[htbp]
\centering
\caption{Exposição gestacional contra marcação por data de nascimento}
\label{tab:exposicao}
\begin{tabular}{lcc}
\hline
Coorte & Fração da gestação pós-ban & Marcação ingênua \\
\hline
dez/2018 & 0{,}00 & 0 \\
jan/2019 & 0{,}00 & 1 \\
mai/2019 & 0{,}44 & 1 \\
out/2019 & 1{,}00 & 1 \\
\hline
\end{tabular}

\vspace{0.5em}
\begin{minipage}{0.9\textwidth}
\footnotesize \emph{Nota:} aritmética de desenho, não resultado. Ao longo de
2019, a marcação por data de nascimento conta doze coortes tratadas onde a
exposição efetiva soma sete --- superestimação de cinco duodécimos.
\end{minipage}
\end{table}

\paragraph{A janela é fixa, e a razão importa.} A retroprojeção usa janela
gestacional \textbf{fixa}, não a duração observada em cada célula. Usar a duração
observada seria endógeno: a prematuridade é um dos desfechos do estudo, e se a
proibição altera a duração da gestação, a janela de exposição se moveria junto
com o tratamento. O desenho passaria a condicionar em variável afetada pelo
tratamento, com viés de direção desconhecida. A gestação observada entra como
desfecho; nunca como definidora de exposição.

\paragraph{Um canal não retroprojeta.} A intoxicação aguda registrada em
internação hospitalar é contemporânea à exposição --- não há gestação entre causa
e efeito. Esse canal usa o mês de internação diretamente. Retroprojetá-lo
inventaria defasagem inexistente.

\subsection{O grupo de comparação}
\label{sub:zero}

O zero produzido pela PAM é zero de \emph{proxy} --- área nula da cultura ---, não
zero de tratamento. Duas fontes de contaminação são conhecidas.

A primeira tem direção. O §2º do art. 28-B, na redação de 2019, proíbe também a
dispersão aérea para controle vetorial. Um município sem agricultura pulverizada,
mas com controle aéreo de endemias, \textbf{é tratado} --- e municípios com
programa de controle vetorial tendem a ser os maiores. A contaminação, portanto,
correlaciona-se com porte populacional, que por sua vez correlaciona-se com
desfecho perinatal. Não é ruído.

A segunda é histórica. O art. 29 da Lei Estadual nº 12.228/1993 autoriza
municípios a legislar supletivamente sobre uso de agrotóxicos, e há registro de
que Limoeiro do Norte tenha proibido a pulverização aérea já em 2009. Se procede,
há unidades \textbf{já tratadas} dentro do grupo de dose alta, e 2015--2018 deixa
de ser período pré-tratamento para todas. O levantamento das leis municipais
anteriores a 2019 é, por isso, anterior em importância ao próprio diagnóstico de
dose.

\paragraph{Quatro construções, todas reportadas.} O grupo de comparação é
construído de quatro maneiras --- área nula da cultura-âncora; área nula de
qualquer cultura candidata a aplicação aérea; ausência de operador aeroagrícola
registrado; e baixa aptidão agroclimática ---, e a dispersão entre elas é
reportada como incerteza sobre o nível.

\begin{quote}
\textbf{Por que a dispersão entre construções é o número, e não uma nota de
rodapé.} O estimador não-paramétrico centra a curva na variação média do grupo de
dose nula. Alterar quem compõe esse grupo não adiciona ruído: \textbf{desloca o
nível da curva inteira}. A escolha do zero é, portanto, aritmética do estimador,
não escolha de apresentação.
\end{quote}

\subsection{Estimação, e três restrições que o desenho declara}
\label{sub:estimacao}

A implementação de referência do estimador impõe restrições que não são detalhe
de software e que, por afetarem o que pode ser afirmado, entram aqui:

\begin{enumerate}
  \item \textbf{O estimador não-paramétrico da curva opera sobre dois períodos.}
    O painel município--mês precisa ser colapsado em pré e pós. Isso é escolha de
    agregação, e está declarada como tal.
  \item \textbf{O \emph{event study} não sai do mesmo estimador que a curva.} Os
    \emph{leads} que dão evidência sobre paralelismo e antecipação vêm de
    estimação separada, com agregação distinta.
  \item \textbf{Covariáveis não são acomodadas.} Qualquer ajuste por
    características municipais ocorre por fora --- na camada de pareamento
    apresentada como robustez ---, não dentro do estimador principal.
\end{enumerate}

\paragraph{A tensão que essas restrições produzem, e que o leitor deve conhecer.}
Colapsar em dois períodos transforma o pré-período em uma média única de
2015--2018. Como essa janela é integralmente posterior ao marco de notícia
(§\ref{sub:dose}), a média pré pode estar contaminada por antecipação --- e, por
não haver \emph{leads} dentro da especificação principal, essa contaminação
\textbf{não é observável nela}. O \emph{event study} existe e é reportado, mas
valida objeto vizinho, não a curva. A tensão é real e não se resolve com
robustez; declara-se, e três mitigações são reportadas: os \emph{leads} como
pré-requisito explícito, a curva com pré-período recuado quando a
comparabilidade da PAM permitir, e a curva com o pré partido ao meio, cuja
diferença mede a contaminação.

\paragraph{Coortes de transição.} O colapso pré/pós não acomoda as coortes cuja
gestação atravessa a vigência. Elas são \textbf{descartadas} --- escolha
conservadora, dado que incluí-las de qualquer lado embute a atenuação da
Tabela~\ref{tab:exposicao} ---, e os cortes são reportados como parâmetro.

\subsection{Poder, e o que conta como resultado negativo}
\label{sub:poder}

O desenho declara seu poder antes de estimar, e não depois.

O parâmetro que decide a detectabilidade não é a dispersão de dose entre
municípios, mas o \textbf{desvio-padrão das tendências municipais do desfecho no
pré-período}. Com poucos municípios de dose alta, é ele que governa o efeito
mínimo detectável. O piso amostral desse desvio --- calculado para município de
porte típico, sem heterogeneidade real, apenas amostragem --- já é da ordem de
17,7 gramas, de modo que cenários mais otimistas devem ser lidos como limite
superior de desempenho, não como caso médio.

\paragraph{O critério de falsificação, declarado antes.} Um resultado nulo é
inconclusivo somente se o intervalo de confiança for largo. Se o intervalo
\textbf{excluir} o efeito que a literatura levaria a esperar, o nulo constitui
evidência de que a proibição não produziu esse efeito --- e o desenho falsifica.
A meia-largura do intervalo é proporcional ao efeito mínimo detectável, de modo
que o critério é aritmético e verificável antes de qualquer estimação.

\begin{quote}
Se o coeficiente estimado ficar \textbf{abaixo} do efeito mínimo detectável do
próprio desenho, ele não constitui achado: constitui \textbf{limite superior
informativo}, e é assim que será reportado.
\end{quote}

\subsection{Inferência}
\label{sub:inferencia}

Com poucos municípios de dose alta, o bootstrap assintótico produz intervalos
estreitos demais. São reportados três procedimentos, lado a lado:

\begin{enumerate}
  \item erro-padrão agrupado convencional, como linha de base;
  \item \emph{wild cluster bootstrap} com pesos de Rademacher e nulo imposto ---
    com a ressalva de que ele sub-rejeita quando os clusters tratados são poucos,
    de modo que um valor-p alto informa mas um valor-p baixo não absolve;
  \item \textbf{inferência por aleatorização} sobre a dose, que não depende de
    aproximação assintótica alguma e é, das três, a mais conservadora.
\end{enumerate}

\begin{quote}
A distância entre os três é informação, não inconveniência. Se o procedimento
convencional indicar significância e a aleatorização não, a leitura correta não é
``o efeito é significativo'': é ``o desenho não distingue''.
\end{quote}

\paragraph{Nota de nomenclatura.} O procedimento de referência para poucas
mudanças de política --- Conley e Taber --- foi construído para tratamento
\textbf{binário}, usando resíduos do grupo de controle como distribuição de
referência. Com tratamento contínuo, a adaptação fiel é permutar a dose. A lógica
é a mesma; o procedimento não. Ele é nomeado aqui pelo que é.

\subsection{Ameaças, e o que cada uma faz ao resultado}
\label{sub:ameacas}

\begin{table}[htbp]
\centering
\caption{Ameaças à identificação e tratamento dado a cada uma}
\label{tab:ameacas}
\small
\begin{tabular}{p{4.2cm}p{4.6cm}p{4.6cm}}
\hline
Ameaça & Efeito sobre o estimando & Tratamento no desenho \\
\hline
Tendências paralelas fortes não verificáveis
& A curva deixa de ser identificada pontualmente
& Limites de identificação parcial sob direção do viés \\

Antecipação a partir de 02/2015
& Contamina a dose, não só o desfecho
& \emph{Leads} até 2018; pré-período recuado; pré partido ao meio \\

Contaminação do grupo de dose nula
& \textbf{Desloca o nível} da curva
& Quatro construções; dispersão reportada como incerteza \\

Municípios já tratados antes de 2019
& \textbf{Falha silenciosa}: resultado sai e está errado
& Levantamento legislativo anterior ao diagnóstico de dose \\

Seleção para nascimento vivo
& Efeito real pode aparecer como nulo
& Série de óbito fetal testada contra a dose antes de qualquer correção \\

Lacuna de fiscalização
& Estimando vira efeito de intenção de tratar
& Pedido de acesso à informação; canal de intoxicação como via independente \\

Poucos clusters tratados
& Precisão aparente maior que a real
& Três procedimentos de inferência reportados juntos \\
\hline
\end{tabular}
\end{table}

\paragraph{Seleção para nascimento vivo, em detalhe.} Se a proibição reduziu
óbito fetal, fetos que antes não chegavam a nascer passam a compor a amostra pela
cauda inferior de peso --- puxando a média na direção \emph{oposta} à do efeito
esperado, e disputando a mesma ordem de grandeza. Antes de qualquer correção, o
desenho verifica se o problema existe: a série de óbito fetal é testada contra a
dose com o mesmo desenho. Se não se move, a discussão encerra-se em limitação
declarada. Se se move, as rotas conhecidas --- reporte conjunto, desfecho
composto, limites de seleção amostral --- são avaliadas, e a escolha é reportada.

\paragraph{Um canal que produz falsificação interna.} A hipótese de substituição
sustenta que, proibido o avião, a aplicação migra para trator e pulverizador
costal, aproximando o veneno do aplicador. Ela prevê movimento nas internações
por intoxicação \textbf{acidental}. Como a proibição alcança o método e não a
molécula, ela prevê \textbf{ausência} de movimento nas intoxicações
\textbf{autoprovocadas}, que respondem à disponibilidade do produto. As duas
séries partilham população, sistema de registro e municípios. Movimento conjunto
não sustenta substituição de método; sustenta algo que move internação em geral.

\subsection{Pré-especificação}
\label{sub:prespec}

O conjunto de desfechos que o desenho torna disponíveis é amplo, e o poder é
curto. Essa combinação --- espaço de busca grande com poder limitado --- é a que
produz achado espúrio com aparência de rigor. Por essa razão, desfecho primário,
exposição primária, conjunto confirmatório e tratamento de multiplicidade são
declarados por escrito antes do primeiro contato com dado real, e o registro é
versionado, de modo que sua anterioridade seja verificável. Todo desvio posterior
é reportado com data e motivo, e o resultado original é preservado.

% 5. Dados e estatísticas descritivas   <- espera a aquisição
% 6. Resultados e robustez              <- FORA DO ALVO DA QUALIFICAÇÃO
% 7. Conclusão

\end{document}
